\subsubsection{Examples of topological spaces}

\begin{example}[Order Topology]
    \phantom{text}\\
    Let $X$ be a set with the order relation $<$. For $a<b$ in $X$, define intervals:
    \[(a,b)=\{x \in X \mid a < x < b\}\]
    \[[a,b)=\{x \in X \mid a \leq x < b\}\]
    Similar for $(a,b]$,$[a,b]$\\
    Let $\beta$ be the set of all intervals in $X$ of the forms:
    \begin{itemize}
        \item $(a,b)$ for $a<b$ in $X$.
        \item $[a_0,b)$ where $a_0$ is the smallest element of $X$.(if it exists).
        \item $(a,b_0]$ where $b_0$ is the largest element of $X$.(if it exists).
    \end{itemize}
    \begin{proposition}
        This is a basis for a topology on $X$.\\
        This is called the order topology on $X$.
    \end{proposition}
    \begin{example}
        Standard topology on $\R$ is order topology, on $\R$
    \end{example}
\end{example}

\begin{example}[Product Topology]
    Let $X,Y$ be topological spaces.\\
    Let $\beta$ be a collection of sets:
    \[\{U \times V \mid U \t{ open in } X\t{, } V \t{ open in } Y\}\]
    \begin{proposition}
        This is a basis for a topology.\\
        This is called the product topology on $X \times Y$
    \end{proposition}
    \begin{example}
        $R \times R$ with product topology is std topology on $\R^2$
    \end{example}
    \phantom{text}\\
    If $\beta_1$ is a basis for topology on $X$, and $\beta_2$ is a basis for topology on $Y$, then 
    \[\{B \times C \mid B \in \beta_1,\,\, C\in \beta_2\}\]
    is a basis for the product topology.
\end{example}
This definition is valid for products of pairs of sets.
\newpage
\begin{lemma}
    Suppose $X$ is a topological space. Let $C$ be a collection of open subsets satisfying:
    \begin{itemize}
        \item for each open set $U \subseteq X$ and each $x \in U$, $\exists c \in C$ s.t. $x \in c \subseteq U$
    \end{itemize}
    Then, $C$ is a basis for top. on $X$.
\end{lemma}
\begin{proof}
    Check $C$ is a basis.
    \begin{itemize}
        \item Every $x \in X$ in some $c \in C$. $X$ is open, so there exists $c$ s.t. $x \in c \subseteq X$
        \item If $c_1,c_2 \in C$, then for any $x \in c_1 \cap c_2$, $\exists c_3 \in C$ s.t. $x \in c_3 \subset c_1 \cap c_2$.\\
        Follows from the fact that $c_1,c_2$ are open which implies $c_1 \cap c_2$ is open so can find $c_3$
    \end{itemize}
    Check: $C$ generates same topology $T$ on $X$.\\
    First, suppose $U$ is open in $T$. By assumption for all $x \in U$, there exists $c \in C$ s.t. $x \in c \subset U$,
    so $U$ is open in the topology generated by $C$.\\
    If $U$ is open in topology generated by $C$. Then $U$ is a union of sets in $C$. But every $c \in C$ is open in $T$, so union $U$ is also open in $T$.
\end{proof}

\begin{example}[Metric Topology]
    \phantom{text}
    \begin{definition}
        Let $X$ be a set. A function $d: X \times X \to \R$ is a metric if it satisfies:
        \begin{itemize}
            \item $d(x,y)=0 \iff x=y$
            \item $d(x,y)=d(y,x) \quad \forall x,y \in X$
            \item $d(x,y)+d(y,z) \geq d(x,z) \quad \forall x,y,z \in X$
        \end{itemize}
    \end{definition}
    A set $X$ equipped with a metric $d$ is a metric space.\\
    Let $\beta$ be the set of all metric balls:
    \[B_d(x,r) = \{y \in X \mid d(x,y)<r\} \quad \t{for $x \in X$, $r>0$}\]
    \begin{proposition}
        This is a basis for a topology\\
        This is called metric topology
    \end{proposition}
    \begin{definition}
        A topological space $X$ is called metrizable if $\exists$metric $d$ on $X$ inducing topology on $X$.
    \end{definition}
    \begin{rmk}
        There could be more than one metric inducing the same topology.
    \end{rmk}
\end{example}